
\section{Discussion and Conclusion}
\label{sec:conclusion}

Our results demonstrate that calibration-free wearable
eye tracking is well suited for collecting gaze data
in real driving. More importantly, they establish a
practical methodology for measuring, improving, and
monitoring gaze quality, providing a foundation for
larger-scale gaze-supervised driving datasets.

\textbf{RQ1: Reliable gaze measurements in real driving.}
Our real-world benchmark demonstrates that Neon achieves
sub-degree median accuracy (0.76--0.83\dg{}) after
person-specific correction using natural driving
references. Only a few fixation events are needed:
four randomly selected scenes already approach the
performance of the pooled calibration set.
Driving-derived references consistently outperform
laboratory-derived offsets, highlighting the value
of evaluating gaze under its intended operating
conditions. Combined with our automated annotation
pipeline, this makes accurate in-situ verification
practical without dedicated calibration procedures.

\textbf{RQ2: Robustness and influencing factors.}
Our extensive controlled experiments largely confirm
Neon's robustness across diverse viewing and
illumination conditions. At the same time, they
reveal meaningful individual differences.
Systematic offsets depend on the wearer--glasses
configuration and can vary with viewing distance
and gaze position, while dynamic tracking primarily
increases dispersion. These observations show why
accuracy should be assessed for the current wearer
and operating conditions rather than assumed from
a general device specification.

\textbf{RQ3: Improving measurement reliability.}
Using these findings, person-specific spatial and
temporal corrections improve upon an already strong
median-offset baseline by 2.5--11.8\% across the
evaluated configurations. Improvements are especially
valuable in dynamic situations, where temporal
filtering reduces larger errors. The results
demonstrate that relatively simple, individually
selected corrections can provide practical benefits
without modifying the underlying gaze estimator.

\textbf{Towards scalable, quality-aware gaze datasets.}
A central contribution of this work is that gaze
data can now be collected together with evidence
about their measurement quality, rather than
treated as error-free supervision labels.
Our pipeline automatically generates reference
observations from natural fixation events, associates
them with vehicle motion, depth, and driving-scenario
annotations, and supports independent target
validation. Complementary quality monitoring
combines semantic viewing-region segmentation,
optical flow, image sharpness, IMU measurements,
and eye-visibility indicators to characterize
sample suitability.

The demonstrated effectiveness of sparse in-situ
references additionally enables lightweight
accuracy checks during acquisition, providing
a mechanism to verify and update wearer-specific
corrections. Together, these components establish
a practical workflow for generating corrected,
quality-annotated gaze observations at scale.


\textbf{Scope and future work.}
The present driving benchmark evaluates two participants
using held-out scenes from recorded drives. Further
recordings across wearers, sessions, and driving conditions
will extend the evaluation of individual variability and
correction transfer. Our automated acquisition pipeline
already combines validated reference generation with
multimodal quality monitoring, providing the tools to
collect larger datasets with spatial references and
sample-level suitability information. Future work will
expand this dataset, investigate data-driven uncertainty
estimation, and make the collected data available to
the research community.


\textbf{Conclusion.}
We have demonstrated that wearable gaze tracking
can provide accurate supervision signals in
real-world driving, and established the experimental
and automated tools to assess their reliability.
By combining validated reference generation,
systematic characterization, sparse correction,
and multimodal quality monitoring, our work moves
towards scalable gaze datasets in which measurement
quality is evaluated rather than assumed.
